\section{Introduction}\label{sec:introduction}

Benchmark contamination, the unintended inclusion of evaluation data in training corpora, poses a growing threat to the integrity of large language model (LLM) evaluation \citep{dodge2021, magar2022, sainz2023}.
As training corpora and public benchmarks both expand, the probability of overlap increases structurally: any model trained on broad web-crawled data may have encountered evaluation items during pretraining or fine-tuning \citep{jacovi2023, chang2024survey}.
When contamination occurs, benchmark scores no longer reflect genuine capability, undermining the primary mechanism by which the community compares and selects models.

A range of detection methods have been proposed, but all require some form of privileged access beyond the binary response record.
N-gram overlap methods need the training corpus itself \citep{dodge2021, magar2022}.
Membership inference approaches such as Min-K\%++ require token-level probabilities \citep{minkprob2024, deng2024}.
Time Travel probing relies on crafted prompts that elicit memorized completions \citep{golchin2024, oren2024}.
Performance-based methods compare results across original and rephrased test variants \citep{constat2024}, and in-context learning approaches require output logits \citep{codec2026}.
These methods form a detection frontier ordered by information access: training data, model weights, logits, token probabilities, rephrased retests, and finally binary patterns.
No existing method operates at the bottom of this frontier.
We ask: \emph{what can be detected from binary response patterns alone?}

The need for minimal-information detection grows more pressing as models evolve.
\citet{wang2025fragility} show that all ten log-probability-based detection methods they evaluate perform at chance level on reasoning models trained with reinforcement learning (RL).
RL fine-tuning reshapes internal token distributions while largely preserving binary correctness patterns, and probability-based detectors cannot distinguish contaminated from clean models under these conditions.
Binary response patterns are structurally invariant to RL-induced distributional changes, as they record what a model answered rather than how its internal distributions shifted.

We observe that a contaminated LLM is functionally analogous to a test-taker with item preknowledge: both produce anomalous response patterns in which difficult items are answered correctly through memorization rather than genuine ability.
This connection points to a mature body of work in psychometrics.
Person-fit statistics, developed over four decades for cheating detection in educational testing \citep{drasgow1985, meijer2001, karabatsos2003, sinharay2017}, quantify exactly this type of response aberrance.
The standardized statistic $\ell_z^*$ \citep{snijders2001} measures how well an individual's response pattern conforms to Item Response Theory (IRT) model predictions, with extreme negative values indicating patterns consistent with preknowledge.
Despite this natural correspondence, to our knowledge no prior work has applied person-fit statistics to LLM benchmark contamination detection.

We introduce IRT-FIT, a framework that applies IRT person-fit analysis at the minimal-information operating point, requiring only binary responses and pre-calibrated item parameters \citep{psnirt2026}.
We compute $\ell_z^*$ for each (model, benchmark) pair and define $\Delta\ell_z^*$ as a difference-in-differences estimator for within-model contamination detection.
In controlled experiments (QLoRA fine-tuning on MMLU), 15\% contamination produces $|\text{Cohen's } d| = 3.42$ ($n{=}3$ clean seeds).
Our closed-form power function relates benchmark size to detection sensitivity ($R^2 = 0.957$), and shows that 80\% power at 15\% contamination requires $n \geq 7{,}958$ items (under MMLU item parameters).
A 4PL versus 3PL ablation reveals a mean $+27.3$ percentage-point power gain on MMLU \citep{hendrycks2021mmlu}, confirming that the upper-asymptote parameter is essential for LLM evaluation contexts.
Across 12 models and 5 benchmarks, $\ell_z^*$ flags 57 of 60 model-benchmark cells for response-pattern anomalies (BH-FDR $\alpha{=}0.05$), versus only 4 flagged by an accuracy $z$-score baseline.
This $14\times$ gap shows that $\ell_z^*$ captures response-pattern structure lost in aggregate accuracy.
We further identify a $\theta$-shift boundary condition: re-estimating ability on contaminated data can degrade sensitivity (non-monotonic dose-response under free $\theta$), but fixing $\theta$ to pre-contamination values restores monotonicity (signal ratio 2.9--3.1 across model families).

Our contributions are:
\begin{enumerate}[leftmargin=*,nosep]
\item We propose IRT-FIT, a person-fit framework for contamination detection at the minimal-information operating point (binary responses only), validated through controlled contamination on Qwen2.5-7B ($|\text{Cohen's } d| = 3.42$ at 15\% contamination, $n{=}3$) and confirmed across model families with Llama-3.1-8B ($|\Delta\ell_z^*| = 3.84$).
\item We derive a closed-form power function ($R^2 = 0.957$) linking benchmark size to detection sensitivity, identify the $\theta$-shift boundary condition, and show that 4PL parameterization yields $+27.3$pp power gain over 3PL.
\item We report three empirical findings from real LLM response data: universal 4PL model misfit (57/60 cells flagged), the inadequacy of accuracy-based anomaly detection (4/60 cells), and a simulation-to-reality gap that motivates the $\Delta\ell_z^*$ design.
\end{enumerate}